Several limitations exist in Orbiter. Here is a list:

\begin{enumerate}
\item
Field elements are encoded as  $\verb'int'$. This limits the size of fields that can be handled to $2^{8s-1}$ where $s=$\verb'sizeof(int)'.
\item
The ranks of elements in the permutation domain are encoded as  $\verb'long int'$. 
This limits the size of permutation domains that can be handled.
The degree of a permutation group must be less that $2^{8s-1}$ where $s=$\verb'sizeof(long int)'.
\item
The finite field class builds tables for the addition and multiplication of field elements. 
This restricts the size of the fields that can be created.
\item
The projective geometry class tries to build a bitmatrix for the adjacency matrix if the number of lines is less than 
\verb'MAX_NUMBER_OF_LINES_FOR_INCIDENCE_MATRIX' which is defined in \verb'src/lib/foundations/geometry/projective_space.cpp'. 
If the number of lines is too big, the table is not created. 
In this case, the projective geometry class may behave slower.
\item
The projective geometry class tries to build a table for the lines if the number of points is less that \verb'MAX_NUMBER_OF_POINTS_FOR_POINT_TABLE' and the number of 
lines is less than 
\verb'MAX_NUMBER_OF_LINES_FOR_LINE_TABLE', 
both of which are  defined in \verb'src/lib/foundations/geometry/projective_space.cpp'.
If the number of points is too big, the table is not created. 
In this case, the projective geometry class may behave slow.
\item
The projective geometry class tries to build a table for the lines through any two points if the number of points is less than 
\verb'MAX_NB_POINTS_FOR_LINE_THROUGH_TWO_POINTS_TABLE' 
which is defined in \verb'src/lib/foundations/geometry/projective_space.cpp'.
If the number of points is too big, the table is not created. 
In this case, the projective geometry class may behave slow.
\item
The projective geometry class tries to build a table for the intersection points of pairs of lines if the number of points is less than 
\verb'MAX_NB_POINTS_FOR_LINE_INTERSECTION_TABLE' 
which is defined in \verb'src/lib/foundations/geometry/projective_space.cpp'.
If the number of points or lines is too big, the table is not created. 
In this case, the projective geometry class may behave slow.
\item
For Windows users: Cygwin by default uses 32 bit integers for both \verb'int' and \verb'long int'. 
Using Cygwin 64 to compile Orbiter recommended.
\item
A limited list of primitive polynomials are hard-coded in Orbiter. 
For large fields, the user must provide their own primitive polynomial.
The polynomials encoded in orbiter are not guaranteed to be compatible with the subfield relationship.
%\item
%The base and strong generators used by Orbiter (see Section~\ref{sec:permutation:groups}) do not lead to optimal Schreier trees. 
%This is a performance issue for projective linear groups. The techniques of shallow Schreier trees (cf. Seress~\cite{Seress03}) 
%would be useful.
%\item
%Orbiter uses long integer arithmetic for the orders of finite groups, so there is no restriction in the order of finite groups from that point of view. 
\end{enumerate}


\bigskip


The command  

\sourcecode{F_15bit}

creates the field $\bbF_{32749}$. Tables are suppressed for speed.
Note that 
$$
2^{15}-19 = 32749
$$
is prime (see~\url{https://primes.utm.edu/lists/2small/0bit.html}).
The command creates two instances of $-1$ in the field and multiplies them. 


\bigskip


The command  

\sourcecode{F_31bit}

creates the field $\bbF_{2147483647}$. Tables are suppressed for speed.
Note that 
$$
2^{31}-1 = 2147483647 = p
$$
is prime (see~\url{https://primes.utm.edu/lists/2small/0bit.html}).
The command creates two instances of $-1$ in the field $\bbZ_p$ 
and multiplies them. 
The fact that this is done correctly is remarkable, because 
the intermediate value exceeds the size of 32 bit signed integers.
Recall that 
$-1$ is stored as $p-1,$ 
so
$$
-1 \equiv 2147483646.
$$
But
$$
-1 \cdot -1 \equiv 2147483646 \cdot 2147483646 = 4611686009837453316,
$$
which exceeds the range of $32$ bit signed integers.
The reason why this command works is that the multiplication of integers 
is performed using long integer arithmetic.
Long integer arithmetic treats integers as text strings, and hence is not bound 
by the size of a machine word. Unfortunately, 
long integer arithmetic is considerably slower than arithmetic in machine words.
So, there is a performance drawback with large finite fields in Orbiter.


\bigskip

The command  

\sourcecode{F_32bit}

creates the field $\bbF_{4294967291}$. 
Note that 
$$
2^{32}-5 = 4294967291
$$
is prime (see~\url{https://primes.utm.edu/lists/2small/0bit.html}).
The command will likely not run, and terminate with an error message. 
The reason is hardware dependent. 
The command tests the size of a machine word.
Most compilers map the C++ data type 
{\em signed int} to 32 bit machine words, 
so this command will fail in such an environment. 
This is true even with 64 bit hardware.
There may be some compiler options to change this.



\bigskip

Machines for which the previous command works correctly, the following command should also be fine:

\sourcecode{F_63bit}

Note that 
$$
2^{63}-25 = 9223372036854775783
$$
is prime (see~\url{https://primes.utm.edu/lists/2small/0bit.html}).
However, the next command will almost surely fail:

\sourcecode{F_64bit}

This is because we are exceeding the range 
that can be stored in a 64 bit unsigned integer machine word.
Note that 
$$
2^{64}-59 = 18446744073709551557
$$
is prime (see~\url{https://primes.utm.edu/lists/2small/0bit.html}).







